\section{A logical simplification: diagonal vacuity only for negations}
\label{sec:negative-diagonal}

This supplement concerns the numerical calculus; it does not alter the
accepted PA194DHS transition table. Replace the full diagonal schema
\[
 \Gamma(B)=B\to U(B,B)
\]
by its negative instances alone,
\[
 \Gamma^-(B)=\Gamma(N(B))
       =N(B)\to U(N(B),N(B)),\qquad N(B)=B\to1.
\]
Keep all six ordinary heads, the dominated implicational head, the
membership-vacuity/pointed-separation gate and its original distinctness
guard, modus ponens, and theorem generalization. Call this calculus
$\mathcal C_-$. All its displayed diagonal instances remain sound under
the inherited PA interpretation, by the same code-freshness lemma.

\subsection{Positive vacuity at the same binder}

For a fixed numerical label $z$, write $\Box A=U(z,A)$ and
$\Diamond A=N(\Box N(A))$. These are abbreviations for the original
numerical constructors. We use the retained heads
\[
 \Diamond\Box A\to A,
 \qquad
 \Box A\to(\Box(A\to B)\to\Box B),
\]
closed generalization, and the derived classical implicational proof
operations. The latter include contraposition, composition, and both
double-negation implications; their derivations use the dominated basis
and the first ordinary head, and use no diagonal or equality-vacuity axiom.

\paragraph{Vacuity-negation lemma.}
From a closed derivation of $A\to\Box A$ one can construct, uniformly,
a closed derivation of $N(A)\to\Box N(A)$.
Indeed contraposition gives
$N(\Box A)\to N(A)$. Generalize this closed theorem and apply the sixth
head to obtain
\[
 \Box N(\Box A)\to\Box N(A).
\]
The first head is $N(\Box N(\Box A))\to A$. Contraposition followed by
double-negation elimination gives
\[
 N(A)\to\Box N(\Box A).
\]
Composition proves the assertion. Generalization is applied only to
closed theorems; there is no assumption-dependent generalization step.
This is the proof operation named \texttt{negationVacuity} in the supplied
\texttt{PA196/Logic/Vacuity.lean}.

\paragraph{Positive-from-negative lemma.}
For arbitrary $B,z$, a closed derivation
\[
                     N(B)\to U(z,N(B))
\]
uniformly yields a closed derivation $B\to U(z,B)$. Apply the preceding
lemma with $A=N(B)$, and compose the three implications
\[
 B\ \longrightarrow\ N(N(B))\
   \longrightarrow\ U(z,N(N(B)))\
   \longrightarrow\ U(z,B).
\]
The first arrow is double-negation introduction. The middle arrow is
vacuity-negation applied to the input proof. The last arrow is obtained
by generalizing the closed double-negation-elimination theorem
$N(N(B))\to B$ and applying universal distribution. Thus replacement
under $U$ is justified by an explicit quantifier-congruence derivation,
not by treating double negations as identical numerical codes.

Applying this lemma to $\Gamma^-(B)$ gives
\begin{equation}
 \mathcal C_-\vdash B\to U(N(B),B).
 \label{eq:negative-fresh-vacuity}
\end{equation}
The binder in (\ref{eq:negative-fresh-vacuity}) is exactly
$z=\operatorname{code}(N(B))$, not $z=\operatorname{code}(B)$.
No relabeling or alpha-conversion of that binder has been asserted.
Since $B<N(B)$, the chosen binder is still larger than every variable
label appearing in $B$.

\subsection{The two places in the inherited PA calibration}

The supplied diagonal calibration replaces primitive equality vacuity in
the original PA196 refutation compiler through two interfaces.

\paragraph{Reverse membership rewriting.}
In that argument the required formula is
$C=E(a,b)\to(\chi(b)\to\chi(a))$, and the separator property is
$\varphi=N(C)$ with separator label $c=\operatorname{code}(\varphi)$.
The needed diagonal clause is exactly
$\varphi\to U(c,\varphi)=\Gamma^-(C)$, which is an axiom of
$\mathcal C_-$. Hence the negative-property specialization, pointed
separation, label-collision cases, and two recovered reverse rewrites
are unchanged. Membership vacuity is retained with its original guard;
it has not been restricted to numerically large binders.

\paragraph{Neighborhood equality.}
For $B=E(a,b)$ choose
\[
                   z_0=\operatorname{code}(N(E(a,b))).
\]
The positive-from-negative lemma gives
$E(a,b)\to U(z_0,E(a,b))$. Lifting the two-way membership rewrites at
$z_0$ then yields
\[
 E(a,b)\to R(a,b,z_0),
\]
where
\[
 R(a,b,z)=U\bigl(z,M(a,z)\leftrightarrow M(b,z)\bigr)
 \wedge U\bigl(z,M(z,a)\leftrightarrow M(z,b)\bigr).
\]
The freshness conditions hold
because $z_0>E(a,b)>a,b$. For any requested $z\ne a,b$, the existing
membership-only binder-transport proof sends $R(a,b,z_0)$ to $R(a,b,z)$;
if $z=z_0$ use the identity proof. That transport uses membership
vacuity, the witness head and the already recovered rewrites, and needs
no primitive equality vacuity. The subsequent neighborhood-witness
interface is obtained by the same existential-monotonicity argument as
before, using this updated neighborhood implication.

These are explicit replacements of the two diagonal-dependent
calibration templates. They suffice to reuse the supplied PA196
refutation compiler at its stated rewriting/neighborhood/witness
interface. Conversely every $\mathcal C_-$ axiom belongs to the original
sound separate-head calculus. Thus, at the same inherited calibration
boundary as the main report, there are primitive-recursive refutation
translations with PA-provable correctness, and
\[
 \mathrm{PA}\vdash
 \left(\exists d\,\operatorname{Prf}_{\mathcal C_-}(d,1)
        \leftrightarrow\neg\operatorname{Con}(\mathrm{PA})\right).
\]
The constructions use fixed finite proof templates, numerical constructor
operations, a decidable binder-equality split and recursion on the input
proof. No uniform arithmetic truth predicate, reflection principle, or
unbounded search for fresh names is required. This does not claim that
$\mathcal C_-$ derives every original $\Gamma(B)$ with its original
binder $B$; the established result is the displayed refutation calibration.

\subsection{Physical cost and the separate positive-offset experiment}

The direct emitter for $\Gamma(N(D))$ uses three copies of the
four-token expression
\texttt{copyD, push, incT, cons}, followed by
\texttt{cons, incT, incT, cons}. It therefore has sixteen tokens in
place of the seven-token ordinary diagonal head. Nine direct insertions
in the fixed A, D and Q grammar layouts, using each of A, D and H as the
source parameter, had optimistic state counts at least 200. These
profiles use the cached-MP, unsymmetric cache-test controller underlying
those layouts, not the accepted DHS controller. These are
compilation-cost screens, not accepted machines.

A separate simple schema emits $\Gamma(D+1)$. It covers all positive
codes, while the omitted $\Gamma(0)$ follows from the truth theorem,
generalization at zero, and propositional weakening. Consequently this
positive-offset variant has a direct proof translation for every
original diagonal instance. The recorded 72-case screen concerned
$\Gamma(R+k)$, with $R\in\{A,D,H\}$, $1\le k\le8$, in the same three
fixed grammar layouts; its minimum optimistic count was 198. Only
$k=1$ has been justified here by the preceding full-schema argument;
other offsets in that screen were exploratory compilation cases.

Both screens subtract a maximum matching of eligible contextual pairs
from the raw control count; such a matching may still fail simultaneous
compatibility. They do not certify all source outcomes, blank startup,
or any global minimum over machines. The data are retained in
\texttt{supplement/negative\_only\_screen.json} and
\texttt{supplement/positive\_offset\_screen.json}. The logical simplification
is proved, but no accepted machine below 194 states follows from it.

The bundled standard-library replay script
\texttt{supplement/semantic\_screens.py} recomputes these 81 profiles and
the 105 earlier controller/radix crosses from three frozen grammar inputs,
using the pinned construction compiler. All 186 recorded profile counts
were reproduced; the controller/radix cross minimum was 194.
